\section*{الفصل \ref{ch:g11:dissolution} --- ذوبان الأجسام الصلبة الأيونية والجزيئية}

\begin{solution}{exo:g11:dissolution:1}
\ce{NaCl(s) -> Na+(aq) + Cl-(aq)}؛
\ce{CaCl2(s) -> Ca^{2+}(aq) + 2Cl-(aq)}؛
\ce{Na2SO4(s) -> 2Na+(aq) + SO4^{2-}(aq)}؛
\ce{AlCl3(s) -> Al^{3+}(aq) + 3Cl-(aq)}.
\end{solution}

\begin{solution}{exo:g11:dissolution:2}
$[\ce{Na+}] = 2 \times 0.050 = \qty{0.10}{mol/L}$؛
$[\ce{SO4^{2-}}] = \qty{0.050}{mol/L}$.
\end{solution}

\begin{solution}{exo:g11:dissolution:3}
\omterm{def:g11:dissolution:solvation}{الإماهة} إحاطة \omterm{def:g9:ions:ion}{الأيون} \omterm{def:g7:atoms-and-molecules:molecule}{بجزيئات} ماء يمسكها به التجاذب.
فالطرف الأكسجيني ($\delta^-$) يتجه نحو \omterm{def:g9:ions:ion}{الكاتيون}، والطرف الهيدروجيني
($\delta^+$) نحو \omterm{def:g9:ions:ion}{الأنيون}: فالشحنات المتعاكسة تتجاذب.
\end{solution}

\begin{solution}{exo:g11:dissolution:4}
$M(\ce{AlCl3}) = 27.0 + 3 \times 35.5 = \qty{133.5}{g/mol}$؛
$n = 2.67 / 133.5 = \qty{0.0200}{mol}$؛ $c = 0.0200 / 0.2000 =
\qty{0.100}{mol/L}$؛ $[\ce{Al^{3+}}] = \qty{0.100}{mol/L}$،
$[\ce{Cl-}] = \qty{0.300}{mol/L}$.
\end{solution}

\begin{solution}{exo:g11:dissolution:5}
الرأس \ce{-COO-} أيوني، يجذبه الماء: \omterm{def:g11:dissolution:hydrophilic}{محب للماء}. والذيل
المكوّن من 17 \omterm{def:g7:atoms-and-molecules:atom}{ذرة} كربون، وليس فيه إلا روابط \ce{C-C} وروابط \ce{C-H} شبه غير
قطبية، غير قطبي: \omterm{def:g11:dissolution:hydrophilic}{كاره للماء}.
\end{solution}

\begin{solution}{exo:g11:dissolution:6}
الملح: الماء (جسم صلب أيوني، \omterm{def:g6:solutions-and-solubility:solute}{ومذيب} قطبي). الشمع: الهكسان الحلقي (سلاسل غير
قطبية). السكر: الماء (فمجموعاته \ce{O-H} تكوّن روابط هيدروجينية مع
الماء). اليود: الهكسان الحلقي (\omterm{def:g11:polarity-and-cohesion:polar-molecule}{جزيء غير قطبي}).
\end{solution}

\begin{solution}{exo:g11:dissolution:7}
مجموعة \ce{O-H} في الإيثانول تكوّن روابط هيدروجينية مع \omterm{def:g7:atoms-and-molecules:molecule}{جزيئات} الماء،
تحل محل التي تكسرها: فيندمج الإيثانول في الماء. أما الهكسان، غير القطبي،
فلا يستطيع تكوين مثل هذه الروابط؛ \omterm{def:g7:atoms-and-molecules:molecule}{وجزيئات} الماء، المشدود بعضها إلى بعض بروابط
هيدروجينية، لا تسمح له بالدخول.
\end{solution}

\begin{solution}{exo:g11:dissolution:8}
تتجنب الذيول الكارهة للماء الماءَ وتذوب في الدهن، في
المركز؛ وتجذب الرؤوسَ المشحونة الماءُ فتبقى في الخارج. وكل
\omterm{def:g11:dissolution:amphiphilic}{المذيلات} تحمل شحنات سالبة على سطوحها، والشحنات المتماثلة
تتنافر: فتبقى متباعدة.
\end{solution}

\begin{solution}{exo:g11:dissolution:9}
الهكسان الحلقي: فهو \omterm{def:g6:solutions-and-solubility:miscible}{غير قابل للامتزاج} بالماء، فيكوّن طبقة منفصلة
يمكن تصريفها، والعطر (غير القطبي) \omterm{def:g2:mixing-and-dissolving:dissolve}{يذوب} فيه جيدًا.
أما الإيثانول \omterm{def:g6:solutions-and-solubility:miscible}{فقابل للامتزاج} بالماء: فلا تتكوّن طبقة ثانية، ولا يمكن فصل
شيء.
\end{solution}

\begin{solution}{exo:g11:dissolution:10}
الهكسان الحلقي أقل كثافة من الماء. وتُصرَّف الطبقة السفلى، المائية،
عبر الصنبور. واليود في الطبقة العليا من الهكسان الحلقي، التي
تبقى في القمع.
\end{solution}

\begin{solution}{exo:g11:dissolution:11}
لا شيء: فكبريتات النحاس أيونية ولا تذوب في الهكسان الحلقي غير
القطبي. فيبقى الماء أزرق والهكسان الحلقي عديم اللون.
\end{solution}

\begin{solution}{exo:g11:dissolution:12}
$[\ce{Cl-}] = 2c$، إذن $c = \qty{0.150}{mol/L}$؛
$n = 0.150 \times 0.2500 = \qty{0.0375}{mol}$؛
$m = 0.0375 \times 111.1 = \qty{4.17}{g}$.
\end{solution}

\begin{solution}{exo:g11:dissolution:13}
ترسّب \omterm{def:g9:ions:ion}{أيونات} الكالسيوم والمغنيسيوم أيوناتِ الستيارات زبدًا:
فيُستهلك الصابون قبل أن يكوّن \omterm{def:g11:dissolution:amphiphilic}{مذيلات}، فيرغي رغوة رديئة.
الكالسيوم وحده: \qty{0.010}{mol} من \ce{Ca^{2+}} تأخذ
\qty{0.020}{mol} من الستيارات، أي $0.020 \times 306.0 =
\qty{6.1}{g}$ من ستيارات الصوديوم في اللتر.
\end{solution}

\begin{solution}{exo:g11:dissolution:14}
الحجم \qty{0.2000}{L}. \ce{Na+}: $0.020 / 0.2000 = \qty{0.10}{mol/L}$؛
\ce{Ca^{2+}}: $0.010 / 0.2000 = \qty{0.050}{mol/L}$؛ \ce{Cl-}:
$(0.020 + 0.020) / 0.2000 = \qty{0.20}{mol/L}$. الشحنة الموجبة:
$0.10 + 2 \times 0.050 = 0.20$؛ والسالبة: $0.20$. متعادل.
\end{solution}

\begin{solution}{exo:g11:dissolution:15}
كتلة لتر من الهبتان \qty{0.68}{kg}، فهو يستطيع أن يحمل
$17 \times 0.68 = \qty{12}{g}$ من اليود، أي نحو 40 مرة أكثر من
لتر من الماء. فإذا رُجّا معًا انتقل اليود كله تقريبًا إلى
الهبتان: \omterm{def:g10:chemical-species:extraction}{فالاستخلاص} فعال.
\end{solution}

\begin{solution}{pb:g11:dissolution:1}
\textbf{1.} من الصخور التي جرى فيها الماء (الحجر الجيري،
والدولوميت، والجبس)، التي تذوب قليلًا.

\textbf{2.} $[\ce{Ca^{2+}}] = 0.120 / 40.1 = \qty{2.99e-3}{mol/L}$؛
$[\ce{Mg^{2+}}] = 0.024 / 24.3 = \qty{9.9e-4}{mol/L}$.

\textbf{3.} \ce{Ca(HCO3)2 -> Ca^{2+}(aq) + 2HCO3-(aq)}؛ إذن
$[\ce{HCO3-}] = 2 \times \num{2.99e-3} = \qty{5.99e-3}{mol/L}$.

\textbf{4.} $\num{2.99e-3} \times 150 = \qty{0.449}{mol}$.

\textbf{5.} $18 \times 12.0 + 35 \times 1.0 + 2 \times 16.0 + 23.0 =
\qty{306.0}{g/mol}$.

\textbf{6.} \ce{C17H35COONa(s) -> C17H35COO-(aq) + Na+(aq)}.

\textbf{7.} الرأس \ce{-COO-}: \omterm{def:g11:dissolution:hydrophilic}{محب للماء}؛ والذيل \ce{C17H35-}:
\omterm{def:g11:dissolution:hydrophilic}{كاره للماء}. وله النوعان من الأجزاء: فهو \omterm{def:g11:dissolution:amphiphilic}{محب للوسطين}.

\textbf{8.} كرية دقيقة من \omterm{def:g9:ions:ion}{أيونات} الصابون، ذيولها نحو الداخل ورؤوسها نحو الخارج.
تذوب الذيول في الدهن، فيتكسر قطيرات تلتف في
\omterm{def:g11:dissolution:amphiphilic}{مذيلات}؛ وتُبقي السطوح المشحونة القطيرات متباعدة في الماء،
الذي يشطفها بعيدًا.

\textbf{9.} يجب أن يكوّن الصابون \omterm{def:g9:ions:ion}{أيونات} \omterm{def:g11:dissolution:amphiphilic}{ومذيلات}، وهذا يتطلب ماءً
حول الرؤوس؛ والدهن يحمله ماء الشطف بعيدًا. أما في
الزيت فيبقى الدهن ذائبًا على الجلد ببساطة.

\textbf{10.} \ce{2C17H35COO- + Ca^{2+} -> Ca(C17H35COO)2}: الشحنات
$2 \times (-1) + 2 = 0$ في اليسار، و$0$ في اليمين.

\textbf{11.} \ce{C17H35COO-}: زيادة $- 2x$؛ \ce{Ca^{2+}}:
$\num{2.99e-3} - x$؛ ستيارات الكالسيوم: $x$. \omterm{def:g9:ions:ion}{فأيونات} الكالسيوم هي
المحدِّدة: $x_{\max} = \qty{2.99e-3}{mol}$.

\textbf{12.} $2 \times \num{2.99e-3} = \qty{5.99e-3}{mol}$ في اللتر.

\textbf{13.} $M = 36 \times 12.0 + 70 \times 1.0 + 4 \times 16.0 + 40.1 =
\qty{606.1}{g/mol}$؛ و$\num{2.99e-3} \times 606.1 = \qty{1.81}{g}$ من الزبد
في اللتر.

\textbf{14.} $2 \times \num{9.9e-4} = \qty{1.98e-3}{mol}$ في اللتر.

\textbf{15.} $2 \times 0.449 = \qty{0.898}{mol}$.

\textbf{16.} $0.898 \times 306.0 = \qty{275}{g}$.

\textbf{17.} المغنيسيوم: $\num{9.9e-4} \times 150 = \qty{0.148}{mol}$،
تأخذ \qty{0.296}{mol} من الستيارات، أي \qty{91}{g} من الصابون؛ وفي المجموع
نحو $275 + 91 = \qty{366}{g}$.

\textbf{18.} $275 / 100 \approx 2.7$ قالب.

\textbf{19.} يضيّع كالسيوم حوض واحد نحو \qty{0.27}{kg} من الصابون
زبدًا.
\end{solution}
